\chapter{Simulation and Tests}
\label{ch:sim}

Each unit has a directed testbench (Section~\ref{sec:directed}), all units share
a test against Berkeley TestFloat (Section~\ref{sec:testfloat}), \sig{fmac} has a
stream test (Section~\ref{sec:streams}) and \sig{fpu} has a mixed-operation test
(Section~\ref{sec:mix}). Section~\ref{sec:coverage} states what the tests cover.
The
commands are in Sections~\ref{sec:make_sim} and~\ref{sec:make_test}.

The flow files support five simulators: Icarus Verilog (the default), xezim,
Vivado xsim, Cadence Xcelium and Verilator 5. They are in \file{lib/sim}, with one
\file{<simulator>.mk} per simulator.

\section{Directed Testbenches}
\label{sec:directed}

Each testbench is \file{<unit>/hw/sim/src/<unit>\_tb.v}. It resets the unit,
applies a fixed sequence of operands and stops. With \sig{VCD=1} it writes the
waveforms to \file{<unit>.vcd}. Table~\ref{tab:tbs} lists what each one applies.

\begin{twotable}{Directed Testbenches}{tab:tbs}{Testbench}{Stimulus and output}{2.6cm}
\sig{fadd\_tb} & Additions of small integers, signed zeros, infinity and subnormal operands, one per clock cycle. It prints nothing; inspect the waveforms. \\ \hline
\sig{fmul\_tb} & Directed multiplications. It prints nothing; inspect the waveforms. \\ \hline
\sig{fdiv\_tb} & Directed divisions. It prints nothing; inspect the waveforms. \\ \hline
\sig{fmac\_tb} & Streams of products: small integers, a simple sum, $\cos(i)$, $\sum 1/i$, $\sum 1/i^2$, a Taylor series of $e^{-2\pi}$, normally distributed values, random values, equally spaced values, heavy cancellation, and one large value with many small ones. It prints each result next to a double-precision reference, with the relative and absolute errors. \\ \hline
\sig{fpu\_tb} & Directed additions (\sig{funct} 0), multiplications (2) and divisions (3). It prints nothing; inspect the waveforms. \\
\end{twotable}

\section{TestFloat}
\label{sec:testfloat}

The test applies vectors from Berkeley TestFloat to each unit, through the
testbench \file{lib/\allowbreak sim/\allowbreak src/\allowbreak testfloat\_tb.v}. The tool \sig{testfloat\_gen}
generates the vectors for the binary16, binary32 and binary64 formats: the smoke
run tests all three, and the regular run binary32. Each
vector holds the operands, the expected result and the expected exception flags.

The test runs every operation in the five rounding modes of
Table~\ref{tab:rnd_mode}: \sig{testfloat\_gen} generates the vectors of each
mode, and the testbench drives the same mode on \sig{rnd\_mode}. The stream and
mixed-operation tests use mode 0.

The testbench applies one vector at a time and waits for \sig{valid\_o}, for at most
1,000 clock cycles. It compares the result bits only, because the units have no
exception flag outputs. Any NaN result matches any expected NaN.

Table~\ref{tab:tf_ops} lists the operations of each unit, and
Table~\ref{tab:tf_runs} the two test runs. For \sig{f32\_mulAdd}, the testbench
sends a stream of two products to \sig{fmac}, $a \cdot b$ and $c \cdot 1.0$, so
the expected result is $a \cdot b + c$.

\begin{twotable}{TestFloat Operations of Each Unit}{tab:tf_ops}{Unit}{Operations}{2.6cm}
\sig{fadd} & \sig{f32\_add} \\ \hline
\sig{fmul} & \sig{f32\_mul} \\ \hline
\sig{fdiv} & \sig{f32\_div} \\ \hline
\sig{fmac} & \sig{f32\_mul}, \sig{f32\_mulAdd} \\ \hline
\sig{fpu} & \sig{f32\_add}, \sig{f32\_sub}, \sig{f32\_mul}, \sig{f32\_div} \\
\end{twotable}

\begin{twotable}{TestFloat Runs}{tab:tf_runs}{Run}{Vectors}{2.6cm}
Smoke & The first 1,000 vectors of each operation. \\ \hline
Regular & TestFloat level 1: 46,464 vectors per two-operand operation. For \sig{fmac}, \sig{f32\_mul} and the first 200,000 of the 6,133,248 level 1 \sig{f32\_mulAdd} vectors. \\ \hline
Full & The smoke run, then the regular run with all 6,133,248 \sig{f32\_mulAdd} vectors for \sig{fmac}. \\
\end{twotable}

The test writes these files to \file{<unit>/hw/sim}:

\begin{itemize}
\item \file{testfloat\_<format>\_<op>\_<mode>.rpt}: the simulator output, ending in a line with the
  number of vectors and failures.
\item \file{fails\_<format>\_<op>\_<mode>.tv}: one line per failing vector, with the operands, the
  expected result and the result of the unit.
\end{itemize}

\section{fmac Stream Test}
\label{sec:streams}

TestFloat applies streams of at most two products to \sig{fmac}. The stream test
applies longer ones. The script \file{fmac/hw/sim/src/fmac\_streams.py} generates
the streams from a fixed seed and checks the results, and the testbench
\file{fmac\_streams\_tb.v} applies them with 2 idle
cycles between streams, and 2 to 4 after every fourth stream.
Table~\ref{tab:streams} lists the stream types. The script uses them in turn, and
every 200th stream is a long one. The smoke run applies 300 streams and the
regular run 6,000.

\begin{twotable}{\sig{fmac} Stream Types}{tab:streams}{Type}{Products}{2.6cm}
Narrow & 1 to 16 products with exponent fields from 120 to 134. \\ \hline
Wide & 1 to 32 products with exponent fields from 60 to 190. \\ \hline
Cancel & 1 to 16 products, their negatives and one smaller product, shuffled. \\ \hline
Ill-conditioned & 8 to 64 products of an ill-conditioned dot product. \\ \hline
Tiny & 1 to 12 products with a zero or subnormal operand. \\ \hline
Special & 1 to 6 products of zeros, infinities, NaN, $\pm 1$ and the smallest subnormals. \\ \hline
Carry & 2 to 64 same-sign products of equal magnitude. \\ \hline
Long & 256, 1,024 or 4,096 products. \\
\end{twotable}

The script compares each result with the exact rational sum of the $n$ products
$a_i b_i$. A stream with a NaN or infinite operand, or whose products all have a
zero operand, must give exactly the IEEE-754 result, with any NaN equal to any
NaN. Any other result passes if it equals
the correctly rounded sum, or if
\[
  |\mathit{result}-\mathit{exact}| \le \frac{\mathrm{ulp}(\mathit{exact})}{2}
  + (n+2)\,2^{-46}\sum_i |a_i b_i| .
\]
An infinite result passes only if it has the sign of the exact sum and the bound
allows the sum to reach the overflow threshold.

The test writes \file{streams.tv}, \file{streams.out}, \file{streams.rpt} and
\file{fails\_streams.tv} to \file{fmac/hw/sim}, and prints the numbers of streams,
correctly rounded results and failures.

\section{fpu Mixed-Operation Test}
\label{sec:mix}

The script \file{fpu/hw/sim/src/fpu\_mix.py} draws a random sequence from the
TestFloat level 1 vectors of \sig{f32\_add}, \sig{f32\_sub}, \sig{f32\_mul} and
\sig{f32\_div}, with weights 30, 30, 35 and 5 and a fixed seed. The testbench
\file{fpu\_mix\_tb.v} issues additions, subtractions and multiplications on
consecutive clock cycles. It holds the operands of a division until its result,
and issues nothing else meanwhile. It checks the results in issue order, with any
NaN equal to any NaN. The smoke run applies 2,000 operations and the regular run
100,000. The test writes \file{mix.tv}, \file{mix.rpt} and \file{fails\_mix.tv} to
\file{fpu/hw/sim}.

\section{Test Coverage}
\label{sec:coverage}

TestFloat covers the binary16 (16/5), binary32 (32/8) and binary64 (64/11)
formats; the other tests of this chapter cover binary32 only. All use the default
values of the other parameters. The model comparison of
Section~\ref{sec:sw_check} also simulates \sig{fpu} and \sig{fmac} at 16/5 and
64/11, and \sig{fmac} with other guard widths.

\sig{make coverage} runs the regular tests in Verilator with line and toggle
coverage, and the regular TestFloat runs also in binary16 and binary64
(\sig{COV\_FMTS}). It merges the runs of each unit and format and writes
\file{<unit>/hw/sim/cov/coverage\_<format>.rpt}, with the coverage of each source
file, the lines never executed and the signals whose bits never toggle both ways,
and annotated sources in \file{<unit>/hw/sim/cov/annotated\_<format>}. A format is
a different circuit, so each has its own report. Every line of every unit is
executed in all three formats. The bits that never toggle are constants by design, such as the low
bits of the special-result encodings, zero extensions and shift limits, and
\sig{rnd\_mode}, which each TestFloat run holds at one mode.
